% Generated draft: requires mathematical and source review.
\documentclass[11pt]{article}
\usepackage[margin=25mm]{geometry}
\usepackage{fontspec}
\setmainfont{Latin Modern Roman}
\usepackage{amsmath,amssymb,mathtools,unicode-math}
\setmathfont{Latin Modern Math}
\usepackage{enumitem,needspace,xurl,hyperref}
\setlength{\parindent}{0pt}
\setlength{\parskip}{5pt}
\setlength{\emergencystretch}{3em}
\newcommand{\N}{\mathbb{N}}
\newcommand{\Z}{\mathbb{Z}}
\newcommand{\Q}{\mathbb{Q}}
\newcommand{\R}{\mathbb{R}}
\newcommand{\C}{\mathbb{C}}
\newcommand{\F}{\mathbb{F}}
\newcommand{\A}{\mathbb{A}}
\newcommand{\PP}{\mathbb{P}}
\newcommand{\E}{\mathbb{E}}
\newcommand{\eps}{\varepsilon}
\newcommand{\dd}{\,\mathrm{d}}
\begin{document}
\section*{UnsolvedMath ID: 10400206}
\textbf{Source problem code:} catalogue entry
\subsection{Definitions and mathematical statement}
For a space \(C\), its integral homology has torsion in degree \(i\) if some nonzero \(x\in H_i(C;\Z)\) satisfies \(mx=0\) for a positive integer \(m\). D. Thurston asks whether the particular configuration spaces used in source reference [237] have such torsion.

The second question asks whether that torsion can yield a torsion invariant of homology 3-spheres. Existence of torsion classes alone does not construct an invariant: one needs a topologically natural rule assigning appropriate torsion data, with independence of choices and the desired nontriviality. The source does not specify whether the target is a fixed torsion group, an obstruction class, or another torsion-sensitive object, and does not assert that every torsion class automatically gives an invariant.

Reference [237] and the configuration-space definitions are absent. Ordinary ordered configurations, unordered configurations, compactifications, relative homology groups and spaces assembled by gluing diagram faces can have different integral homology; none is substituted for the named construction. The configuration size, homological degree, dependence on the manifold, and intended homology-sphere category also need that reference. In the usual integral interpretation a homology 3-sphere has the integer homology of \(S^3\), but this convention requires source confirmation. No finite-type degree or coefficient prime is specified.
\subsection{Short English statement}
Do the specific configuration spaces in reference [237] have integral homology torsion, and can it produce a torsion-sensitive invariant of homology 3-spheres? The spaces and the intended invariant construction are not supplied.
\subsection{Sources}
\begin{itemize}
\item Local UnsolvedMath registry, permanent ID 10400206, source code catalogue entry; statement attribution D. Thurston; proposer field and proposal year absent; update 2026-08-06.
\item Ohtsuki, \emph{Problems on invariants of knots and 3-manifolds} (2002), Problem 11.10, PDF page 157: \url{https://msp.org/gtm/2002/04/gtm-2002-04-024s.pdf}. Registry access date 2026-07-29; PDF font-structure extraction. Provenance distinguishes 231 numbered bold headings. Present primary-literature status and redistribution terms require review.
\item The crucial source reference [237] is named without its bibliography entry. The supplied summary gives no additional precise configuration-space definition or verified result.
\end{itemize}
\end{document}
